\exercisesection{Commutative Banach Algebras}

In the exercises below, all algebras considered are over C, unless explicitly stated otherwise.

\begin{enumerate}
\def\labelenumi{\arabic{enumi})}
\item
  Let A be a commutative Banach algebra and I a maximal ideal of A. Show that there are only two possible cases: either I is the kernel of a character of A, or I is a hyperplane containing \(A^{2}\) . (Use exerc. 6 of I, p.~154). In particular, if I is not closed, I is a dense hyperplane of A containing \(A^{2}\) . (To obtain an example of this latter situation, consider exerc. 3.)
\item
  Let A be the Banach algebra of sequences \(x = (x_{n})_{n \geqslant 1}\) of complex numbers tending to 0, with \(\|x\| = \sup_{i} |x_{i}|\) . Let I be the set of elements of A with finite support. Then I is a dense ideal of A, and is contained in no maximal ideal. (Observe that \(A^{2} = A\) , and use exerc. 1.)
\item
  Let \(\Lambda\) a set, \(p \in [1, +\infty[\) , and \(\mathbf{A} = \ell^p(\Lambda)\) the set of families \(x = (x_{\lambda})_{\lambda \in \Lambda}\) , where \(x_{\lambda} \in \mathbf{C}\) and
\end{enumerate}

\[
\| x \| _ {p} = \left(\sum_ {\lambda} | x _ {\lambda} | ^ {p}\right) ^ {1 / p} <   + \infty .
\]

\begin{enumerate}
\def\labelenumi{\alph{enumi})}
\item
  For componentwise addition and multiplication, and the norm \(x \mapsto \| x \|_p\) , the space A is a Banach algebra.
\item
  We have \(A^{2} \neq A\) and \(A^{2}\) is dense in A.
\item
  The space \(\mathsf{X}(\mathsf{A})\) of nonzero characters of A is naturally identified with \(\Lambda\) endowed with the discrete topology. (Use the fact that the dual of \(\ell^{p}(\Lambda)\) is \(\ell^{q}(\Lambda)\) , where q is the conjugate exponent of p.) The Gelfand transformation then becomes the identity map and its image in the space of continuous functions tending to 0 at infinity on \(\Lambda\) is dense and not closed.
\end{enumerate}

\begin{enumerate}
\def\labelenumi{\arabic{enumi})}
\setcounter{enumi}{3}
\tightlist
\item
  Let \((\alpha_{n})_{n\geqslant 0}\) a sequence in \(\mathbf{R}_{+}^{*}\) such that \(\alpha_0 = 1\) , \(\alpha_{m + n}\leqslant \alpha_m\alpha_n\) for all integers \(m\) and \(n\) , and such that \(\alpha_{n}^{1 / n}\to 0\) . Let A be the set of formal series \(x = \sum_{n = 0}^{\infty}x_{n}z^{n}\in \mathbf{C}[[z]]\) such that
\end{enumerate}

\[
\| x \| = \sum_ {n = 0} ^ {\infty} | x _ {n} | \alpha_ {n} <   + \infty .
\]

Show that A is a commutative unital Banach algebra topologically generated by z, and that the unique maximal ideal of A is the set of \(x \in A\) without constant term. (Observe that z is quasi-nilpotent.)

\begin{enumerate}
\def\labelenumi{\arabic{enumi})}
\setcounter{enumi}{4}
\item
  In the algebra \(\mathbf{C}(\mathrm{X})\) of rational fractions over \(\mathbf{C}\) , the ideal \(\{0\}\) is maximal, but is not of codimension 1.
\item
  Let \(\Delta\) the closed unit disk in C. Show that the group of invertible elements in \(\mathcal{C}(\Delta)\) is not dense in \(\mathcal{C}(\Delta)\) .
\item
  Let A be a commutative unital Banach algebra and \(\chi \in \mathsf{X}(\mathsf{A})\) . For every \(x \in \mathsf{A}\) , set \(\|x\|_{\chi} = |\chi(x)| + \|x - \chi(x)\|\) . Show that \(x \mapsto \|x\|_{\chi}\) is a norm, that \(\|xy\|_{\chi} \leqslant \|x\|_{\chi}\|y\|_{\chi}\) , \(\|e\|_{\chi} = 1\) , and that if \(\|e\| = 1\) , then \(\|x\| \leqslant \|x\|_{\chi} \leqslant 3\|x\|\) .
\item
  Let A be a commutative unital Banach algebra such that \(\|1\| = 1\) . Let N be the set of norms equivalent to the given norm, for which A is still a Banach algebra, and for which 1 has norm 1.
\end{enumerate}

\begin{enumerate}
\def\labelenumi{\alph{enumi})}
\tightlist
\item
  Let \(x_{0} \in A\) such that \(\varrho(x_{0}) < 1\) . For every \(x \in A\) , we set:
\end{enumerate}

\[
\left\| x \right\| ^ {\prime} = \inf \sum_ {n} \left\| a _ {n} \right\|
\]

for all representations of \(x\) in the form \(a_0 + a_1 x_0 + \cdots + a_n x_0^n\) (with \(a_0, a_1, \ldots, a_n \in \mathrm{A}\) ). Show that \(x \mapsto \|x\|^{\prime}\) is an element of N. (Since \(\varrho(x_0) < 1\) , there exists \(k\) such that \(\|x_0^n\| \leqslant k\) for all \(n\) , whence \(\|x\| \leqslant k \|x\|^{\prime}\) .)

\begin{enumerate}
\def\labelenumi{\alph{enumi})}
\setcounter{enumi}{1}
\item
  Show that \(\|x_{0}\|^{\prime}\leqslant1\) .
\item
  Deduce from a) and b) that, for every \(x \in A\) , one has
\end{enumerate}

\[
\varrho (x) = \inf _ {n \in \mathrm{N}} n (x).
\]

\begin{enumerate}
\def\labelenumi{\alph{enumi})}
\setcounter{enumi}{3}
\tightlist
\item
  If the only element of N is the given norm, we have \(A = C \cdot 1\) . (Use c) and exerc. 7.)
\end{enumerate}

\begin{enumerate}
\def\labelenumi{\arabic{enumi})}
\setcounter{enumi}{8}
\item
  Let X be a locally compact space and A the Banach algebra \(\mathcal{C}_{0}(X)\) . Let I be an ideal of A such that, for every \(x \in X\) , there exists \(f \in I\) such that \(f(x) \neq 0\) . Show that I contains the ideal of functions with compact support in X.
\item
  Let A be a commutative Banach algebra and D a continuous derivation of A into A, that is, a continuous linear map of A into A such that \(\mathrm{D}(ab)=(\mathrm{Da})b+a(\mathrm{Db})\) for \(a,b\in A\) .
\end{enumerate}

\begin{enumerate}
\def\labelenumi{\alph{enumi})}
\tightlist
\item
  For every \(\chi \in \mathsf{X}'(\mathbf{A})\) and every \(\lambda \in \mathbf{C}\) , the series
\end{enumerate}

\[
\varphi_ {\lambda} (a) = \sum_ {n = 0} ^ {\infty} \frac {\lambda^ {n}}{n !} \chi (\mathrm{D} ^ {n} (a))
\]

converges; \(\lambda \mapsto \varphi_{\lambda}(a)\) is an entire function, and \(a \mapsto \varphi_{\lambda}(a)\) is a character of A.

\begin{enumerate}
\def\labelenumi{\alph{enumi})}
\setcounter{enumi}{1}
\item
  The number \(\varphi_{\lambda}(a)\) is independent of \(\lambda\) . (Observe that \(|\varphi_{\lambda}(a)| \leqslant \|a\|\) according to \(a\) ).) Hence \(\chi(\mathrm{D}(a)) = 0\) .
\item
  The image of D is contained in the radical of A. Deduce from this a new proof of assertion f) of exerc. 31 of I, p.~162.
\end{enumerate}

\begin{enumerate}
\def\labelenumi{\arabic{enumi})}
\setcounter{enumi}{10}
\tightlist
\item
  Let B be the algebra over C of infinitely differentiable functions on {[}0,1{]} in C.
\end{enumerate}

\begin{enumerate}
\def\labelenumi{\alph{enumi})}
\tightlist
\item
  Let A be a subalgebra of B, equipped with a norm making A a Banach algebra. Show that there exists a sequence \((m_{n})_{n\geqslant0}\) into \(R_{+}\) such that, for all \(x\in A\) , one has
\end{enumerate}

\[
\sup _ {0 \leqslant t \leqslant 1} | x ^ {(n)} (t) | = \mathrm{O} (m _ {n}).
\]

(Let \(\mathrm{B}_n\) the Banach algebra of functions \(n\) times continuously differentiable on {[}0,1{]}. By prop. 9 of I, p.~40, the canonical injection \(\mathrm{A} \to \mathrm{B}_n\) is continuous; let \(\mathrm{M}_n\) be its norm; if \(x \in \mathrm{A}\) , one has \(\sup_{0 \leq t \leq 1} |x^{(n)}(t)| \leqslant n!\mathrm{M}_n \| x\|\) .)

\begin{enumerate}
\def\labelenumi{\alph{enumi})}
\setcounter{enumi}{1}
\tightlist
\item
  There exists no norm on B making it a Banach algebra.
\end{enumerate}

\begin{enumerate}
\def\labelenumi{\arabic{enumi})}
\setcounter{enumi}{11}
\tightlist
\item
  Let A be a unital Banach algebra and \((x_{n})\) a sequence of invertible elements of A tending to an element \(x \in A\) If the sequence \((\varrho(x_{n}^{-1}))\) is bounded, and if \(x_{n}x = xx_{n}\) for all n, then x is invertible. (By the cor. of prop. 8 of I, p.~38, \(\varrho(1 - x_{n}^{-1}x) \leqslant \varrho(x_{n}^{-1})\varrho(x_{n} - x)\) , hence \(x_{n}^{-1}x\) is invertible for n sufficiently large.)
\end{enumerate}

\begin{enumerate}
\def\labelenumi{¶ \arabic{enumi})}
\setcounter{enumi}{12}
\tightlist
\item
  Let A be the Banach algebra \(\ell^{2}(\mathbf{N})\) (exerc. 3). Let \(\mathrm{A}_{0}\) the subalgebra formed by sequences with finite support. Let B be the direct sum of \(\mathrm{A}_{0}\) and of \(\mathbf{C}\) , equipped with the multiplication \((f,\alpha)(g,\beta)=(fg,0)\) for \(f,g\in\mathrm{A}_{0},\alpha,\beta\in\mathbf{C}\) , and the norm
\end{enumerate}

\[
\left\| (f, \alpha) \right\| = \sup \left(\left\| f \right\|, \left| \alpha - \sum_ {n} f (n) \right|\right).
\]

\begin{enumerate}
\def\labelenumi{\alph{enumi})}
\item
  The completed algebra \(\widehat{\mathbf{B}}\) is a commutative Banach algebra whose radical R is \(\mathbf{C} \cdot (0,1)\) , and \(\widehat{\mathbf{B}}/\mathbf{R}\) is isometrically isomorphic to A.
\item
  Let \(\pi\colon\widehat{\mathrm{B}}\to\widehat{\mathrm{B}}/\mathrm{R}\) the canonical morphism. There exists no subalgebra \(\mathrm{B}_{1}\) of \(\widehat{\mathrm{B}}\) supplementary to R such that \(\pi|\mathrm{B}_{1}\colon\mathrm{B}_{1}\to\widehat{\mathrm{B}}/\mathrm{R}=\mathrm{A}\) is a homeomorphism. (Let \(\mathrm{B}_{1}\) be such a subalgebra. Let \(u_{n}\in\mathrm{A}\) such that \(u_{n}(n)=1\) , \(u_{n}(m)=0\) for \(m\neq n\) . Let \(e_{n}\in\mathrm{B}_{1}\) such that \(\pi(e_{n})=u_{n}\) . Show that \(e_{n}\in\mathrm{A}_{0}\) , that \(\sum_{n=1}^{\infty}\frac{1}{n}u_{n}\) converges in A but \(\sum_{n=1}^{\infty}\frac{1}{n}e_{n}\) does not converge in \(\widehat{\mathrm{B}}\) .)
\end{enumerate}

\begin{enumerate}
\def\labelenumi{\arabic{enumi})}
\setcounter{enumi}{13}
\tightlist
\item
  In \(\mathbf{C}^2\) , let \(\Omega\) the compact set defined by
\end{enumerate}

\[
\frac {1}{2} \leqslant | z _ {1} | ^ {2} + | z _ {2} | ^ {2} \leqslant 1,
\]

and let A be the subalgebra of \(\mathcal{C}(\Omega)\) consisting of functions that are holomorphic in \(\mathring{\Omega}\) . Let U be the open set defined by \(|z_1|^2 + |z_2|^2 < 1\) . By virtue of a theorem of Hartogs (cf.~J.-P. DEMAILLY, Complex analytic and differential geometry, Ch. I, th. 3.28), for every \(f \in \mathrm{A}\) , there exists a function \(\widetilde{f} \in \mathcal{C}(\overline{\mathrm{U}})\)

, unique, which is holomorphic in U and coincides with f on \(\Omega\) . Show that the canonical injection from A into \(\mathcal{C}(\Omega)\) is an isometric morphism h whose image is a full subalgebra of \(\mathcal{C}(\Omega)\) but \(\mathsf{X}(h)\) is not surjective.

\begin{enumerate}
\def\labelenumi{\arabic{enumi})}
\setcounter{enumi}{14}
\item
  Let A and B be commutative unital Banach algebras, \(\varphi\) a unital morphism from A to B and \((x_{\lambda})_{\lambda \in \Lambda}\) a family of elements of A such that the closed full subalgebra of A generated by the \(x_{\lambda}\) is equal to A. For \(\mathsf{X}(\varphi)\) to be surjective, it is necessary and sufficient that \(\mathrm{Sp}_{\mathrm{A}}^{\Lambda}((x_{\lambda})) = \mathrm{Sp}_{\mathrm{B}}^{\Lambda}((\varphi (x_{\lambda})))\) (Use diagram (1) of no. 6 of I, p.~41, where the right-hand arrow is bijective by prop. 12 of I, p.~43.)
\item
  Let A be a commutative Banach algebra. The following conditions are equivalent:
\end{enumerate}

\begin{enumerate}
\def\labelenumi{(\roman{enumi})}
\item
  A is radical-free and \(\mathcal{G}(\mathrm{A})\) is closed in \(\mathcal{C}_0(\mathsf{X}(\mathrm{A}))\) ;
\item
  \(\mathcal{G}\) is a homeomorphism from A onto \(\mathcal{G}(\mathrm{A})\) ;
\item
  there exists a constant c \textgreater{} 0 such that \(\|x\|^{2} \leqslant c\|x^{2}\|\) for all \(x \in A\) .
\end{enumerate}

\begin{enumerate}
\def\labelenumi{\arabic{enumi})}
\setcounter{enumi}{16}
\item
  Let A be a commutative Banach algebra and I a closed ideal of A. Denote by h the canonical injection of I into A. The space \(\mathsf{X}'(\mathsf{I})\) is identified with the quotient space of \(\mathsf{X}'(\mathsf{A})\) by the equivalence relation defined by \(\mathsf{X}'(h)\) . (Use part (a) of Exercise 4 of I, p.~153).
\item
  Let A be a unital commutative Banach algebra and \(x\) an element of A. Suppose that the closed full subalgebra of A generated by \(x\) is equal to A, so that the map \(\chi \mapsto \chi(x)\) allows one to identify \(\mathsf{X}(\mathsf{A})\) and \(\mathrm{Sp}_{\mathsf{A}}(x)\) . Let \(\mathcal{H}\) the set of functions \(\mathcal{G}y\) on \(\mathrm{Sp}_{\mathsf{A}}(x)\) , for \(y\) ranging over A. Then \(\check{\mathrm{S}}_{\mathcal{H}}(\mathrm{Sp}_{\mathsf{A}}(x))\) (INT, IV, §7, no. 4) is the boundary of \(\mathrm{Sp}_{\mathsf{A}}(x)\) relative to C. (To show that \(\check{\mathrm{S}}_{\mathcal{H}}(\mathrm{Sp}_{\mathsf{A}}(x))\) is contained in this boundary, use the maximum principle. Conversely, let \(z_0\) a point of this boundary; to show that \(z_0 \in \check{\mathrm{S}}_{\mathcal{H}}(\mathrm{Sp}_{\mathsf{A}}(x))\) , consider \((x - z_1)^{-1}\) for \(z_1\) quite close to \(z_0\) into \(\mathbf{C} = \mathrm{Sp}_{\mathsf{A}}(x)\) .)
\item
  Let A be a commutative unital Banach algebra, B a closed unital subalgebra of A, and T the map \(\chi \mapsto \chi|B\) from X(A) into X(B), and let R be the equivalence relation on X(A) defined by T. Then T induces, by passing to the quotient, a homeomorphism from X(A)/R onto T(X(A)), and for every \(x \in B\) , the function \(f = \mathcal{G}_{\mathrm{B}}(x)\) is such that \(|f|\) attains its upper bound on T(X(A)).
\item
  Let \(\Delta\) the disk \(|z| \leqslant 1\) into \(\mathbf{C}\) the set of \(f \in \mathcal{C}(\Delta)\) that are holomorphic in the interior of \(\Delta\) and \(\mathrm{A}_1\) the Banach subalgebra of \(\mathcal{C}(\Delta)\) generated by A and by the function \(z \mapsto |z|\) . Show that \(\mathrm{X}(\mathrm{A}_1)\) is homeomorphic to the set of \((x_{1}, x_{2}, x_{3}) \in \mathbf{R}^{3}\) such that \(x_{1}^{2} + x_{2}^{2} \leqslant x_{3}^{2}\) and \(0 \leqslant x_{3} \leqslant 1\) .
\item
  \begin{enumerate}
  \def\labelenumii{\alph{enumii})}
  \tightlist
  \item
    Let \((\mathrm{X}_{\lambda})_{\lambda\in\Lambda}\) a family of indeterminates. For any formal series \(f\in\mathbf{C}[[(\mathrm{X}_{\lambda})]]_{\lambda\in\Lambda}\) , let \(\|f\|\) the sum of the absolute values of its coefficients. Let A be the set of \(f\in\mathbf{C}[[(\mathrm{X}_{\lambda})]]_{\lambda\in\Lambda}\) such that \(\|f\|<+\infty\) Equipped with the usual addition and multiplication, A is a commutative unital Banach algebra without radical.
  \end{enumerate}
\end{enumerate}

\begin{enumerate}
\def\labelenumi{\alph{enumi})}
\setcounter{enumi}{1}
\tightlist
\item
  For every commutative unital Banach algebra B, there exists an algebra A of the type in a) and a continuous unital morphism from A onto B.
\end{enumerate}

\begin{enumerate}
\def\labelenumi{\arabic{enumi})}
\setcounter{enumi}{21}
\tightlist
\item
  Let \(\Omega\) a compact subset of \(\mathbf{C}^n\) . The map
\end{enumerate}

\[
(z _ {1}, \ldots , z _ {n}) \longmapsto (z _ {1}, \ldots , z _ {n}, \overline {{z}} _ {1}, \ldots , \overline {{z}} _ {n})
\]

is a homeomorphism from \(\Omega\) onto a compact polynomially convex subset of \(\mathbf{C}^{2n}\) . (Let \(A = \mathcal{C}(\Omega)\) , \(z_i\) the coordinate functions on \(\mathbf{C}^n\) . Then \((z_i|\Omega, \overline{z}_i|\Omega)_{1 \leqslant i \leqslant n}\) is a topological generating system of A, and the map in question is the map from \(X(A) = \Omega\) into \(\mathbf{C}^{2n}\) defined by this topological generating system.)

\begin{enumerate}
\def\labelenumi{\arabic{enumi})}
\setcounter{enumi}{22}
\tightlist
\item
  We denote \(\Omega\) the set of \((z_1, z_2) \in \mathbf{C}^2\) such that \(|z_1| \leqslant 1\) , \(|z_2| \leqslant 1\) . Let \(\alpha \in [0,1]\) and \(\Omega_\alpha\) the set of \((z_1, z_2) \in \mathbf{C}^2\) such that
\end{enumerate}

\[
| z _ {1} | \leqslant 1, \quad | z _ {2} | \leqslant (1 - \alpha) | z _ {1} | + \alpha .
\]

\begin{enumerate}
\def\labelenumi{\alph{enumi})}
\item
  The complements of \(\Omega\) and \(\Omega_{\alpha}\) into \(\mathbf{C}^2\) are connected.
\item
  \(\Omega\) is the polynomially convex hull of \(\Omega_{\alpha}\) . (If \((z_1^0, z_2^0) \in \Omega\) and if \(\mathrm{P} \in \mathbf{C}[z_1, z_2]\) , there exists \(z \in \mathbf{C}\) such that \(|z| = 1\) and \(|\mathrm{P}(z_1^0, z_2^0)| \leqslant |\mathrm{P}(z, z_2^0)|\) ; but \(|\mathrm{P}(z, z_2^0)| \leqslant \sup_{(z_1, z_2) \in \Omega_{\alpha}} |\mathrm{P}(z_1, z_2)|\) .)
\end{enumerate}

\begin{enumerate}
\def\labelenumi{\arabic{enumi})}
\setcounter{enumi}{23}
\item
  Let \(n \geqslant 1\) . There exists a closed convex subset of \(\mathbf{C}^{n}\) that is not polynomially convex.
\item
  Let A be a complete commutative unital locally m-convex algebra (cf.~I, p.~165, exerc. 37).
\end{enumerate}

\begin{enumerate}
\def\labelenumi{\alph{enumi})}
\item
  An element x of A is invertible if and only if \(\chi(x) \neq 0\) for every continuous character \(\chi\) of A. The radical of A is the intersection of the kernels of the continuous characters of A, hence is closed.
\item
  In example e) of exercise 37 of I, p.~165, let I be the ideal formed by the f such that \(f(n) = 0\) for every integer n \textgreater{} 0 sufficiently large. Then every maximal ideal containing I is dense and of infinite codimension.
\end{enumerate}

\begin{enumerate}
\def\labelenumi{\arabic{enumi})}
\setcounter{enumi}{25}
\tightlist
\item
  Let X be a compact topological space and E a complex Banach space, whose norm is denoted by \(\|\cdot\|_{E}\) . We denote by \(\mathcal{C}(X,E)\) the complex vector space of continuous functions from X into E.
\end{enumerate}

\begin{enumerate}
\def\labelenumi{\alph{enumi})}
\tightlist
\item
  Equipped with the norm
\end{enumerate}

\[
\| f \| = \sup _ {x \in \mathrm{X}} \| f (x) \| _ {\mathrm{E}},
\]

the space \(\mathcal{C}(\mathrm{X},\mathrm{E})\) is a Banach space.

\begin{enumerate}
\def\labelenumi{\alph{enumi})}
\setcounter{enumi}{1}
\tightlist
\item
  Let B be the vector subspace of \(\mathcal{C}(\mathrm{X},\mathrm{E})\) generated by the functions \(x\mapsto f(x)y\) , where \(f\in \mathcal{C}(\mathrm{X})\) and \(y\in \mathrm{E}\) . Show that B is dense in \(\mathcal{C}(\mathrm{X},\mathrm{E})\) .
\end{enumerate}

In what follows, we consider a commutative Banach algebra A.

\begin{enumerate}
\def\labelenumi{\alph{enumi})}
\setcounter{enumi}{2}
\tightlist
\item
  Equipped with the multiplication \((fg)(x) = f(x)g(x)\) for all \(x\in \mathrm{X}\) , the Banach space \(\mathcal{C}(\mathrm{X},\mathrm{A})\) is a commutative Banach algebra.
\end{enumerate}

\begin{enumerate}
\def\labelenumi{\alph{enumi})}
\setcounter{enumi}{3}
\tightlist
\item
  Every character \(\chi\in\mathsf{X}(\mathcal{C}(\mathrm{X},\mathrm{A}))\) is of the form \(f\mapsto\widetilde{\chi}(f(x))\) , where \(\widetilde{\chi}\in\mathsf{X}(\mathrm{A})\) and \(x\in\mathrm{X}\) . (First treat the case where A is unital and \(\|1\|=1\) ; consider the restriction of \(\chi\) to \(\mathcal{C}(\mathrm{X})\cdot1\) and to the image of the homomorphism \(\mathrm{A}\to\mathcal{C}(\mathrm{X},\mathrm{A})\) which sends \(a\) to the constant function \(x\mapsto a\) ).
\end{enumerate}

\begin{enumerate}
\def\labelenumi{\alph{enumi})}
\setcounter{enumi}{4}
\tightlist
\item
  The map from \(\mathsf{X}(\mathrm{A})\times \mathrm{X}\) into \(\mathsf{X}(\mathcal{C}(\mathrm{X},\mathrm{A}))\) which to \((\widetilde{\chi},x)\) associates the character \(f\mapsto \widetilde{\chi} (f(x))\) is a homeomorphism.
\end{enumerate}

\begin{enumerate}
\def\labelenumi{\arabic{enumi})}
\setcounter{enumi}{26}
\tightlist
\item
  Let A be a commutative unital Banach algebra over C. A closed subset F of X(A) is called a boundary of A if one has
\end{enumerate}

\[
\sup _ {\chi \in \mathsf {X} (\mathrm{A})} | \chi (x) | = \sup _ {\chi \in \mathrm{F}} | \chi (x) |
\]

for all \(x \in A\) .

\begin{enumerate}
\def\labelenumi{\alph{enumi})}
\item
  Order the set of boundaries of A by inclusion. There exists a minimal boundary.
\item
  Let S be a minimal boundary of A. Every boundary of A contains S, hence S is the intersection of all boundaries of A. (Let F be a boundary of A; show that for every \(x \in S\) , every neighborhood U of x meets F; for this, use the fact that S --- U is not a boundary of A.)
\end{enumerate}

One says that S is the Shilov boundary of A, and one denotes it by \(\partial A\) .

\begin{enumerate}
\def\labelenumi{\alph{enumi})}
\setcounter{enumi}{2}
\item
  If the image of the Gelfand transform of A is dense in \(\mathcal{C}(\mathrm{X}(\mathrm{A}))\) , the Shilov boundary of A is equal to \(\mathrm{X}(\mathrm{A})\) .
\item
  Let A be the algebra of continuous functions on the unit disk \(\Delta\) of \(z \in C\) such that \(|z| \leqslant 1\) that are holomorphic inside of \(\Delta\) (example 9 of I, p.~20). Then \(\mathsf{X}(\mathsf{A})\) is identified with \(\Delta\) (exercise 6 of I, p.~193) and the Shilov boundary of A is identified with the unit circle U.
\end{enumerate}

\begin{enumerate}
\def\labelenumi{\arabic{enumi})}
\setcounter{enumi}{27}
\item
  Let X be a compact topological space and n an integer \(\geqslant 1\) . Let I be a two-sided ideal of the Banach algebra \(\mathcal{C}(\mathrm{X},\mathrm{M}_{n}(\mathbf{C}))\) . Show that there exists a unique closed subset S ⊂ X such that I is the set of \(f \in \mathcal{C}(\mathrm{X}, \mathrm{M}_{n}(\mathbf{C}))\) whose restriction to S is zero.
\item
  * Let X be a compact real manifold. We equip \(\mathcal{C}^{\infty}(\mathrm{X})\) of the topology of uniform convergence of functions and all their derivatives; it is a Fréchet space. Let A be a subalgebra of \(\mathcal{C}(\mathrm{X})\) We assume that A is equipped with a norm \(f \mapsto \| f\|_{\mathrm{A}}\) which makes it a Banach algebra, such that the canonical injection of A into \(\mathcal{C}(\mathrm{X})\) is continuous. We further assume that \(\mathcal{C}^{\infty}(\mathrm{X})\) is contained and dense in A.
\end{enumerate}

\begin{enumerate}
\def\labelenumi{\alph{enumi})}
\tightlist
\item
  There exists a constant \(C \geqslant 0\) and an integer \(k \in N\) such that for all \(f \in \mathcal{C}^{\infty}(X)\) , one has
\end{enumerate}

\[
\| f \| _ {\mathrm{A}} \leqslant \operatorname * {C} \sup _ {x \in \mathrm{X}} | f ^ {(k)} (x) |.
\]

\begin{enumerate}
\def\labelenumi{\alph{enumi})}
\setcounter{enumi}{1}
\item
  The subalgebra A is full in \(\mathcal{C}(\mathrm{X})\) . (Let \(f \in \mathrm{A}\) that does not vanish on X and \(\delta = \inf_{x \in \mathrm{X}} |f(x)|\) ; choose \(g \in \mathcal{C}^{\infty}(\mathrm{X})\) such that \(\|f - g\|_{\mathrm{A}} < \delta/3\) , then note that \(f^{-1} = g^{-1} \sum_{n \in \mathbf{N}} ((g - f)/g)^{n}\) into \(\mathcal{C}(\mathrm{X})\) and verify that the series converges in A.)
\item
  Deduce from this a new proof of Wiener's theorem (I, p.~38, example).
\end{enumerate}

